\begin{algorithm}[H]
\SetAlgoLined
\setstretch{0.92}
\SetAlFnt{\normalsize}
\SetArgSty{textnormal}
\SetKwInput{KwIn}{输入}
\SetKwInput{KwInit}{初始化}
\SetKwInput{KwOut}{输出}
\SetKwFor{For}{for}{do}{end for}
\SetKwFor{While}{while}{do}{end while}
\SetKwIF{If}{ElseIf}{Else}{if}{then}{else if}{else}{end if}
\SetKw{KwRet}{return}
\SetKwFunction{SolveLinear}{SolveLinear}
\caption{问题二微小时间段拓展分离变量法}
\label{alg:q2-thermal}
\KwIn{三层材料参数$\rho_i,c_i,k_i,d_i\;(i=1,2,3)$，
DSC校正积分面积$A_{\rm DSC}$，
扫描速率$v$，
人体侧换热系数$h_b=4.0\ \mathrm{W/(m^2\cdot K)}$，
外侧对流系数$h_e=20.96\ \mathrm{W/(m^2\cdot K)}$，
环境温度$T_\infty=-40\,^{\circ}\mathrm C$，
时间步$\Delta t=0.25\ \mathrm s$，
模态数$M=40$，
总时长$t_{\max}$}
\KwInit{计算潜热$L_{\rm pcm}=60A_{\rm DSC}/v$
及权重函数$w(T)$\;
预求第一层Robin--Dirichlet本征值
$\tan(\lambda_{1,m}d_1)=-k_1\lambda_{1,m}/h_b$\;
预求第三层Dirichlet--Robin本征值
$\tan(\lambda_{3,m}d_3)=-k_3\lambda_{3,m}/h_e$\;
设$T_i(x,0)=37\,^{\circ}\mathrm C$,\;
$n\leftarrow0$,\;$t_{15}\leftarrow\varnothing$,\;$t_{10}\leftarrow\varnothing$}
\While{$n\Delta t < t_{\max}$
    且$T_1(0,n\Delta t)>10\,^{\circ}\mathrm C$}{
    由当前温度场计算中间层平均温度$\bar T_2^n$
    及有效比热$c_{2,\rm eff}^n=c_2+L_{\rm pcm}w(\bar T_2^n)$\;
    更新第二层热扩散率
    $\alpha_2^n=k_2/(\rho_2 c_{2,\rm eff}^n)$\;
    假设当前段内界面温度线性变化，
    对每层构造模态展开\;
    组装两处材料界面热流连续条件，
    形成$2\times2$线性方程组\;
    $[\beta_{12}^n,\beta_{23}^n]^\top
     \leftarrow\SolveLinear(\cdot)$\;
    由模态展开推进各层温度场至$t_{n+1}$\;
    \If{相变固化进度已达$1.0$}{
        冻结该单元有效比热为$c_2$，
        不再继续释放潜热\;
    }
    \If{$t_{15}=\varnothing$ \rm 且 $T_1(0,t_{n+1})\le15\,^{\circ}\mathrm C$}{
        $t_{15}\leftarrow t_{n+1}$\;
    }
    \If{$t_{10}=\varnothing$ \rm 且 $T_1(0,t_{n+1})\le10\,^{\circ}\mathrm C$}{
        $t_{10}\leftarrow t_{n+1}$\;
    }
    $n\leftarrow n+1$\;
}
\KwRet $t_{15}$、$t_{10}$及各时刻温度场$T_i(x,t)$\;
\KwOut{贴身侧温度首次降至$15\,^{\circ}\mathrm C$的时间$t_{15}$，
降至$10\,^{\circ}\mathrm C$的时间$t_{10}$}
\end{algorithm}